\section{Introduction}
\label{sec:intro}

Can large language models (LLMs) generate \emph{good} scientific ideas? Skepticism
is widespread: model-generated hypotheses tend to be either incremental or
implausibly radical, and ``novelty'' itself resists definition~\citep{si2024llm}.
A recurring answer in recent work is to \emph{ground} ideation rather than prompt
for ideas abstractly. SciMON grounds in retrieved literature~\citep{wang2023scimon};
HypoGeniC grounds in labeled data~\citep{zhou2024hypogenic}; ResearchAgent grounds
in structured proposals~\citep{baek2024researchagent}; and SGA grounds in simulation
feedback~\citep{ma2024sga}. Each reports that its grounded method beats an
ungrounded prompt.

\para{Why this matters.} If we can say \emph{which kinds of grounding} turn raw LLM
ideation into useful ideas---and at what cost---we give practitioners a concrete
recipe and move the field past single-number ``novelty'' leaderboards toward
grounded usefulness. This is increasingly urgent because the ungrounded baseline is
a moving target: the 2023-era finding that grounding helps was measured against much
weaker models, and a strong modern LLM may already encode much of what grounding was
meant to supply.

\para{The gap.} Every prior method grounds in \emph{one} way and compares to an
ungrounded prompt. Nobody has held task, model, generation budget, and evaluation
metric fixed and ablated the grounding \emph{mechanisms} against each other. So we
cannot say whether literature comparison, induction, empirical feedback, deduction,
or proposal writing is responsible for the reported gains---or whether, with a strong
2025 model, any of them still helps. We also heed the field's central caveat, that
novelty is over-trusted and can flip once ideas are executed~\citep{si2024llm}, by
scoring \emph{grounded usefulness} rather than a bare novelty score.

\para{Our approach.} We build the first single-protocol, head-to-head ablation of
six ideation regimes: one ungrounded control (C0) and five grounding mechanisms
(deduction, induction, empirical small-experiments/outcome-imagination, literature
comparison, and proposal writing). We operationalize an ``idea'' as a \emph{task
hypothesis}---an explicit decision rule explaining a binary classification
phenomenon---and ``better'' as how accurately an LLM applying a bank of $15$ such
hypotheses predicts held-out labels, in-distribution (\IND) and out-of-distribution
(\OOD). Holding the model (\gptmini), bank size, and an identical evaluator fixed,
the \emph{only} variable is how the ideas were grounded. We run all six regimes on
three \hypogenic{} tasks across three seeds ($54$ banks, $\approx 11.9$k API calls,
$\approx\$3$), and complement accuracy with diversity, a novelty--usefulness
analysis, and a cost/benefit accounting.

\para{What we find.} Grounding does not uniformly help. Pooled across regimes,
grounded banks are statistically indistinguishable from---and slightly below---the
ungrounded control (\IND{} $-1.2$ pts, \OOD{} $-2.5$ pts, both n.s.). But the
mechanism matters decisively: empirical small-experiment grounding is the only
regime that significantly improves \IND{} accuracy ($+4.4$ pts, Cohen's $d=1.09$,
$p=0.02$), while proposal writing significantly \emph{hurts} ($-5.7$ pts, $d=-1.56$,
$p=0.008$) and collapses diversity. Data grounding can \emph{overfit}: on deceptive
reviews the ungrounded model transfers near-perfectly \OOD{} ($0.97$) using general
priors, whereas data-grounded rules latch onto dataset-specific surface cues and drop
to $0.78$--$0.82$ \OOD. Empirical grounding's gains cost roughly $80\times$ the
generation tokens of the control.

\para{Contributions.}
\begin{itemize}[leftmargin=*,itemsep=0pt,topsep=2pt]
    \item We conduct the first single-protocol, head-to-head ablation of six LLM
    ideation regimes, holding model, task suite, generation budget, and evaluation
    fixed and scoring each by objective downstream utility (\IND{} and \OOD{}
    accuracy).
    \item We show that grounding does not uniformly help a strong LLM, and isolate
    the one mechanism that does: grounding in empirical feedback (test-and-revise)
    is the only regime that significantly improves in-distribution usefulness.
    \item We surface two cautions the field under-reports: pure elaboration
    (proposal writing) actively degrades idea quality and diversity, and data
    grounding can trade out-of-distribution robustness for in-distribution fit.
    \item We complement accuracy with diversity, a novelty--usefulness analysis (a
    task-difficulty confound that vanishes within task), and a cost/benefit
    accounting, and release all banks, predictions, and statistics for reproducibility.
\end{itemize}

The rest of the paper reviews related work (\secref{sec:related}), details the
protocol and six regimes (\secref{sec:method}), presents results
(\secref{sec:results}), and discusses implications and limitations
(\secref{sec:discussion}).
